% Optional strengthening; do not use naive weakening for this calculus.
\section{A further restriction: only the newest parameter is needed}

In this section numerical parameter contexts are finite ordered lists.
Every eigenparameter is put at the front of its list. Let
$\mathcal C_0$ have the five rules of $\mathcal C_V$, except that its
VSub source must be the front parameter of the current numerical
context. In de Bruijn notation that source is $v_0$. This variant
requires a different metaproof: primitive-tag-preserving weakening of
an arbitrary $\mathcal C_0$ derivation is false.

\begin{definition}
A derivation with initial numerical parameter list $\Xi$ is
\emph{internally sourced} if, at every VSub node, its source parameter
was introduced by an E eigenparameter binder strictly inside that
derivation, rather than belonging to the initial list $\Xi$.
Moreover, that E binder is the most recently opened numerical binder
still active at the VSub node; there is no intervening active E or I
eigenparameter binder.
Thus the ordinary newest-source side condition holds at every
VSub node; the definition additionally records the source's origin.
\end{definition}

\begin{lemma}[Safe structural changes for internally sourced proofs]
An internally sourced $\mathcal C_0$ derivation admits injective
renaming of its initial numerical parameters into any larger initial
parameter list, and weakening by clauses using only those initial
parameters. The result is internally sourced relative to the enlarged
initial list.
\end{lemma}
\begin{proof}
First freshen the finitely many eigenparameters away from all inserted
data. Preserve every internal eigenparameter and its binder, and apply
the given injection only to the marked initial parameters. Traverse
the proof tree while retaining those original marks: the initial list
is not reset when descending into a premise. Equivalently, carry
the number and names of active internal binders as an induction
parameter. An I or E extension puts its unchanged fresh parameter ahead of
all pre-existing parameters. At every VSub node the source is an
internal parameter. Its origin in the most recently opened active
numerical E binder, and hence its newest-source status, cannot be
changed by inserting, permuting or renaming the initial parameters,
all of which remain older than that source. Capture-free substitution
commutes with injective renaming, normalization is equivariant, and
renaming preserves the literal-zero mode test. Clause insertion does
not change numerical order. All other checks are unchanged.

Equivalently, in de Bruijn notation, initial-parameter renaming acts
only on indices at or above the number of locally active eigenbinders.
At any VSub node that number is positive and the source is index zero,
so the source index is unchanged. This is not a proof of unrestricted
weakening for arbitrary $\mathcal C_0$ derivations.
\end{proof}

\begin{theorem}
Every $\mathcal C_T$ derivation has an internally sourced
$\mathcal C_0$ translation with the identical end sequent and initial
numerical parameter list.
\end{theorem}
\begin{proof}
Induct on the given $\mathcal C_T$ proof tree. Recursively translated
premises are internally sourced relative to their own initial
parameter lists. Copy Refl, Use, E and I nodes unchanged. Combining
the translated premises at these nodes preserves the invariant:
a source that was introduced inside a premise is also introduced
inside the combined derivation.

For a crossed-transport node, apply the direct joint-name construction
from the preceding section to the translated premises. Whenever that
construction carries a premise proof into a new numerical scope or
adds a local clause, apply the safe structural lemma to that entire
premise proof. This is legitimate because it is internally sourced by
the induction hypothesis. The equality-free numerical-case macro
uses only Refl, Use, E and I in addition to these safely weakened
premises. Thus no VSub source is created by any of its structural
operations.

The construction's one new VSub is at the end of the branch of the
fresh-name E binder. Its source $b$ is the newest numerical parameter
of that branch. The major and minor premises are obtained from the
new clause directly by Use and Refl, so no intervening eigenbinder
is active there. Hence the VSub meets the source-zero restriction.
Its source was introduced by that E binder inside the constructed
proof. No E or I eigenparameter is introduced between that binder
and the final VSub node. Thus the entire replacement is internally
sourced relative to the original initial list. Close E and complete the induction.
\end{proof}

The reverse translation from $\mathcal C_0$ to $\mathcal C_T$ is the
same local translation of VSub into crossed transport as before.
Therefore $\mathcal C_0$ and $\mathcal C_T$ prove exactly the same
sequents. General weakening, arbitrary numerical substitution, and
unrestricted crossed transport for arbitrary $\mathcal C_0$ premise
proofs are now admissible by transferring those proofs to
$\mathcal C_T$ and translating the resulting proof back. There is no
circular use of those structural facts in the completeness proof.

The same invariant and translation work for the native six-rule
presentation: replace Use by QE, E by DE, and copy assumption and
reflexivity. The direct fresh-name construction's final VSub still
occurs under its own newest DE eigenparameter.
